\section{Part II: The Branch Point}

\subsection{Locus}

This section establishes a critical distinction between isotropic and anisotropic scale patterns. When the scale pattern is fully isotropic (distinguishing no direction), all rotational labels vanish, and no defects exist. Conversely, any nonzero rotational structure forces the pattern to be anisotropic. A seventh axiom of observability (A7) makes $k = 3$ a theorem rather than an assumption, and ensures that no scale-invariant fundamental interactions exist in the framework.

\begin{definition}[\lean{Locus.Observability}]
A dimension parameter $m$ satisfies the \textbf{Observability} condition if the rotational dimension equals the scale dimension: that is,
\[
\text{rotDim}_2(m) = \text{scaleDim}_2(m).
\]
\end{definition}

\begin{theorem}[\lean{Locus.observability\_eq\_postulate}]
For any $m \in \mathbb{N}$, the Observability condition at $m$ is equivalent to the Observability postulate $\text{A7}(m)$.
\begin{proof}
This is immediate from the definitions; the two formulations are identical.
\end{proof}
\end{theorem}

\begin{theorem}[\lean{Locus.three\_from\_observability}]
If $m \in \mathbb{N}$ satisfies the Observability condition, then $m = 2$.
\begin{proof}
By definition, Observability means $\text{rotDim}_2(m) = \text{scaleDim}_2(m)$. The theorem that rotational and scale dimensions match exactly when $m = 2$ (from the dimensional analysis in \texttt{Dimension.lean}) yields the result.
\end{proof}
\end{theorem}

\begin{theorem}[\lean{Locus.observability\_fails\_elsewhere}]
For any $m \in \mathbb{N}$ with $m \neq 2$, the Observability condition fails at $m$.
\begin{proof}
Suppose the Observability condition held. Then by \texttt{three\_from\_observability}, we would have $m = 2$, contradicting $m \neq 2$. Thus the Observability condition cannot hold.
\end{proof}
\end{theorem}

\begin{theorem}[\lean{Locus.observability\_iff\_three}]
For any $m \in \mathbb{N}$, the Observability condition holds at $m$ if and only if $m = 2$.
\begin{proof}
Observability at $m$ is equivalent to $\text{rotDim}_2(m) = \text{scaleDim}_2(m)$, and dimensional analysis shows this occurs precisely when $m = 2$.
\end{proof}
\end{theorem}

\begin{theorem}[\lean{Locus.no\_scale\_invariant\_interaction}]
Let $u_0, \kappa, \rho, t_0, t, t' \in \mathbb{R}$ with $t \neq t'$ and $\kappa \rho \neq 0$. Then the inverse-square couplings at different scales are unequal:
\[
\text{invSqCoupling}(u_0, \kappa, \rho, t_0, t) \neq \text{invSqCoupling}(u_0, \kappa, \rho, t_0, t').
\]
\begin{proof}
Suppose for contradiction that the couplings are equal. Then their definitions give
\[
u_0 + \kappa(\rho(t - t_0)) = u_0 + \kappa(\rho(t' - t_0)),
\]
which simplifies to $\kappa \rho t = \kappa \rho t'$. Since $\kappa \rho \neq 0$, cancellation yields $t = t'$, contradicting $t \neq t'$.
\end{proof}
\end{theorem}

\begin{theorem}[\lean{Locus.constant\_coupling\_iff\_no\_content}]
Let $u_0, \kappa, \rho, t_0, t \in \mathbb{R}$ with $\kappa \rho = 0$. Then the inverse-square coupling at any scale is constant:
\[
\text{invSqCoupling}(u_0, \kappa, \rho, t_0, t) = u_0.
\]
\begin{proof}
By definition, $\text{invSqCoupling}(u_0, \kappa, \rho, t_0, t) = u_0 + \kappa(\rho(t - t_0))$. Since $\kappa \rho = 0$, the product $\kappa \rho (t - t_0) = 0$, so the coupling equals $u_0$.
\end{proof}
\end{theorem}

\begin{definition}[\lean{Locus.Isotropic}]
A pattern $v \in \text{Fin}(k) \to \mathbb{R}$ (an assignment of $k$ real values) is \textbf{isotropic} if it is constant, i.e., $v_i = v_j$ for all indices $i, j \in \text{Fin}(k)$.
\end{definition}

\begin{theorem}[\lean{Locus.isotropic\_pair\_parallel}]
Let $v, w \in \text{Fin}(k) \to \mathbb{R}$ be two isotropic patterns. Then for any indices $i, j \in \text{Fin}(k)$,
\[
v_i w_j = v_j w_i.
\]
\begin{proof}
Since $v$ is isotropic, $v_i = v_j$. Since $w$ is isotropic, $w_i = w_j$. Therefore $v_i w_j = v_j w_j = v_j w_i$.
\end{proof}
\end{theorem}

\begin{theorem}[\lean{Locus.isotropic\_no\_rotation}]
Let $v, w \in \text{Fin}(k) \to \mathbb{R}$ be two isotropic patterns. Then for any indices $i, j \in \text{Fin}(k)$, the wedge (exterior) product vanishes:
\[
\text{wedge}(v, w)_{ij} := v_i w_j - v_j w_i = 0.
\]
\begin{proof}
By the isotropic pair parallel theorem, $v_i w_j = v_j w_i$ for all $i, j$. Thus $\text{wedge}(v, w)_{ij} = v_i w_j - v_j w_i = 0$.
\end{proof}
\end{theorem}

\begin{theorem}[\lean{Locus.labels\_require\_anisotropy}]
Let $v, w \in \text{Fin}(k) \to \mathbb{R}$ and $i, j \in \text{Fin}(k)$. If the wedge product at $(i,j)$ is nonzero,
\[
\text{wedge}(v, w)_{ij} \neq 0,
\]
then at least one of the patterns is anisotropic: $\neg \text{Isotropic}(v) \lor \neg \text{Isotropic}(w)$.
\begin{proof}
Suppose for contradiction that both $v$ and $w$ are isotropic. Then by the isotropic no rotation theorem, $\text{wedge}(v, w)_{ij} = 0$ for all $i, j$, contradicting the assumption that $\text{wedge}(v, w)_{ij} \neq 0$.
\end{proof}
\end{theorem}

\begin{theorem}[\lean{Locus.everything\_switches\_on\_together}]
Let $v \in \text{Fin}(k) \to \mathbb{R}$ be an isotropic pattern. Then the following hold simultaneously:
\begin{enumerate}
\item For all $i, j \in \text{Fin}(k)$, $\text{wedge}(v, v)_{ij} = 0$ (no rotational label).
\item For all $i, j \in \text{Fin}(k)$, $v_i = v_j$ (no direction is distinguished).
\end{enumerate}
\begin{proof}
Both statements follow directly from the definition of isotropy. The first is the isotropic no rotation theorem applied to $v$ with itself; the second is the definition of isotropy.
\end{proof}
\end{theorem}

\begin{theorem}[\lean{Locus.particle\_locus\_is\_anisotropic}]
Let $v, w \in \text{Fin}(k) \to \mathbb{R}$ and $i, j \in \text{Fin}(k)$. If the wedge product is nonzero,
\[
\text{wedge}(v, w)_{ij} \neq 0,
\]
then the pair $(v, w)$ cannot both be isotropic: $\neg(\text{Isotropic}(v) \land \text{Isotropic}(w))$.
\begin{proof}
Suppose both $v$ and $w$ are isotropic. Then by the isotropic no rotation theorem, $\text{wedge}(v, w)_{ij} = 0$, contradicting the hypothesis. Therefore the pair cannot both be isotropic.
\end{proof}
\end{theorem}

\subsection{Dynamics}

This section translates the coupling structure between the scale and rotation sectors into consequences for the dynamics and spectrum of the theory. The evolution equation is shown to be Bianchi identity, which forces conservation of energy-momentum. Within multiplets (degeneracy blocks of the directional scale), curvature vanishes because parallel transport operators commute. Curvature is generated only between multiplets, and multiplet structure changes only where the scale pattern changes.

\begin{theorem}[\lean{Dynamics.no\_flatness\_from\_commuting\_derivations}]
Let $A \in \text{Fin}(n) \to M$ be a sequence of elements in a differential ring where all partial derivatives vanish: $D_i(x) = 0$ for all $i$ and all $x \in M$. Then for any $i, j \in \text{Fin}(n)$, the curvature is purely commutator-valued:
\[
F(A)_{ij} = \text{ad}(A_i)(A_j) := [A_i, A_j].
\]
\begin{proof}
The curvature is defined as $F(A)_{ij} = D_i(A_j) - D_j(A_i) - [A_i, A_j]$. Since all derivatives vanish, $D_i(A_j) = 0$ and $D_j(A_i) = 0$, leaving $F(A)_{ij} = [A_i, A_j]$.
\end{proof}
\end{theorem}

\begin{theorem}[\lean{Dynamics.nonflat\_witness}]
The two directional scaling operators $K_1, K_2$ from the coupling satisfy:
\begin{enumerate}
\item $\text{ad}(K_1)(K_2) = J$.
\item $J \in \text{Matrix}(\text{Fin}(3), \text{Fin}(3), \mathbb{R})$ is nonzero.
\end{enumerate}
\begin{proof}
The bracket $[K_1, K_2]$ equals the rotation $J$ by the definition of the coupling's boost-rotation relation. The matrix $J$ is nonzero because its $(1,2)$-entry is nonzero.
\end{proof}
\end{theorem}

\begin{theorem}[\lean{Dynamics.commutative\_sector\_transports\_commute}]
Let $C$ be a commutative ring with a differential ring structure of order $n$, and let $A \in \text{Fin}(n) \to C$. For any $i, j \in \text{Fin}(n)$ and $m \in C$, the covariant derivatives in the commutative sector commute:
\[
\text{covD}_A^{(i)}(\text{covD}_A^{(j)}(m)) - \text{covD}_A^{(j)}(\text{covD}_A^{(i)}(m)) = 0.
\]
\begin{proof}
The commutator of covariant derivatives is $[\text{covD}_A^{(i)}, \text{covD}_A^{(j)}](m) = \text{ad}(F(A)_{ij})(m)$. In a commutative ring, $\text{ad}(x)(m) = [x, m] = xm - mx = 0$ for all $x, m$ by commutativity. Thus the covariant derivatives commute.
\end{proof}
\end{theorem}

\begin{theorem}[\lean{Dynamics.commutative\_curvature\_undetectable}]
Let $C$ be a commutative ring with a differential ring structure of order $n$, and let $A \in \text{Fin}(n) \to C$. For any $i, j \in \text{Fin}(n)$ and $m \in C$, the curvature acts trivially:
\[
\text{ad}(F(A)_{ij})(m) = 0.
\]
\begin{proof}
In a commutative ring, the adjoint action $\text{ad}(x)(m) = [x, m] = xm - mx = 0$ by commutativity.
\end{proof}
\end{theorem}

\begin{theorem}[\lean{Dynamics.source\_conserved\_of\_field\_equation}]
Let $A \in \text{Fin}(n) \to M$, $\kappa \in M$, and $T \in \text{Fin}(n) \to \text{Fin}(n) \to M$. If the curvature is proportional to a source tensor,
\[
F(A)_{ij} = \kappa \cdot T_{ij} \quad \text{for all } i, j,
\]
then the source satisfies the cyclic covariant divergence condition for all $i, j, k$:
\[
\text{covD}_A^{(i)}(\kappa \cdot T_{jk}) + \text{covD}_A^{(j)}(\kappa \cdot T_{ki}) + \text{covD}_A^{(k)}(\kappa \cdot T_{ij}) = 0.
\]
\begin{proof}
Substitute the field equation into the left-hand side:
\[
\text{covD}_A^{(i)}(F(A)_{jk}) + \text{covD}_A^{(j)}(F(A)_{ki}) + \text{covD}_A^{(k)}(F(A)_{ij}).
\]
The Bianchi identity (from \texttt{Connection.lean}) states that the cyclic covariant divergence of the curvature vanishes identically, so this sum equals zero.
\end{proof}
\end{theorem}

\begin{theorem}[\lean{Dynamics.no\_field\_equation\_of\_nonconserved}]
Let $A \in \text{Fin}(n) \to M$, $\kappa \in M$, and $T \in \text{Fin}(n) \to \text{Fin}(n) \to M$. If the cyclic divergence of $\kappa T$ fails to vanish at some point,
\[
\text{covD}_A^{(i)}(\kappa \cdot T_{jk}) + \text{covD}_A^{(j)}(\kappa \cdot T_{ki}) + \text{covD}_A^{(k)}(\kappa \cdot T_{ij}) \neq 0,
\]
then $\kappa T$ cannot be the source of any connection's curvature. That is, there does not exist a connection with
\[
F(A)_{ij} = \kappa \cdot T_{ij} \quad \text{for all } i, j.
\]
\begin{proof}
Suppose such a field equation held. Then by the source conserved of field equation theorem, the cyclic divergence would vanish, contradicting the hypothesis. Therefore no such field equation can exist.
\end{proof}
\end{theorem}

\begin{theorem}[\lean{Dynamics.source\_antisymm}]
Let $A \in \text{Fin}(n) \to M$, $\kappa \in M$, and $T \in \text{Fin}(n) \to \text{Fin}(n) \to M$. If the curvature is proportional to the source,
\[
F(A)_{ij} = \kappa \cdot T_{ij} \quad \text{for all } i, j,
\]
then the source inherits the curvature's antisymmetry: for all $i, j$,
\[
\kappa \cdot T_{ij} = -(\kappa \cdot T_{ji}).
\]
\begin{proof}
The curvature is antisymmetric: $F(A)_{ij} = -F(A)_{ji}$ by definition. Substituting the field equation yields $\kappa \cdot T_{ij} = -(\kappa \cdot T_{ji})$.
\end{proof}
\end{theorem}

\begin{theorem}[\lean{Dynamics.field\_equation\_content}]
Let $A \in \text{Fin}(n) \to M$, $\kappa \in M$, and $T \in \text{Fin}(n) \to \text{Fin}(n) \to M$. If the curvature is proportional to the source,
\[
F(A)_{ij} = \kappa \cdot T_{ij} \quad \text{for all } i, j,
\]
then both conservation and antisymmetry hold for all $i, j, k$:
\begin{enumerate}
\item $\text{covD}_A^{(i)}(\kappa \cdot T_{jk}) + \text{covD}_A^{(j)}(\kappa \cdot T_{ki}) + \text{covD}_A^{(k)}(\kappa \cdot T_{ij}) = 0$ (cyclic divergence vanishes).
\item $\kappa \cdot T_{ij} = -(\kappa \cdot T_{ji})$ (antisymmetry).
\end{enumerate}
\begin{proof}
Both statements follow directly from the source conserved of field equation and source antisymm theorems.
\end{proof}
\end{theorem}

\begin{definition}[\lean{Dynamics.Realizes}]
Let $s \in \text{Fin}(N) \to A^\times$ be an assignment of units to directions, and $S \in \text{Fin}(N) \to \text{Matrix}(m, m, \mathbb{R})$ be an assignment of scaling operators. We say $S$ \textbf{realizes} $s$ if equal units have equal scaling operators: for all $a, b \in \text{Fin}(N)$,
\[
s_a = s_b \implies S_a = S_b.
\]
\end{definition}

\begin{theorem}[\lean{Dynamics.same\_multiplet\_no\_rotation}]
Let $s \in \text{Fin}(N) \to A^\times$ and $S \in \text{Fin}(N) \to \text{Matrix}(m, m, \mathbb{R})$, where $A$ is a commutative ring. Suppose $S$ realizes $s$. If two directions $a, b \in \text{Fin}(N)$ are in the same multiplet (i.e., in the same orbit under the action determined by equal scale values),
\[
\text{SameOrbit}_s(a, b),
\]
then the bracket of their scaling operators vanishes:
\[
\text{br}(S_a, S_b) = [S_a, S_b] = 0.
\]
\begin{proof}
By definition of SameOrbit, $s_a = s_b$. Since $S$ realizes $s$, we have $S_a = S_b$. A scaling operator brackets to zero with itself by the coupling structure, so $[S_a, S_b] = [S_a, S_a] = 0$.
\end{proof}
\end{theorem}

\begin{theorem}[\lean{Dynamics.rotation\_implies\_different\_multiplet}]
Let $s \in \text{Fin}(N) \to A^\times$ and $S \in \text{Fin}(N) \to \text{Matrix}(m, m, \mathbb{R})$, where $A$ is a commutative ring. Suppose $S$ realizes $s$. If the bracket of scaling operators at two directions is nonzero,
\[
\text{br}(S_a, S_b) \neq 0,
\]
then the directions cannot be in the same multiplet:
\[
\neg \text{SameOrbit}_s(a, b).
\]
\begin{proof}
Suppose for contradiction that $\text{SameOrbit}_s(a, b)$. Then by the same multiplet no rotation theorem, $\text{br}(S_a, S_b) = 0$, contradicting the hypothesis. Therefore the directions must be in different multiplets.
\end{proof}
\end{theorem}

\begin{theorem}[\lean{Dynamics.larger\_multiplet\_fewer\_generators}]
Let $s \in \text{Fin}(N) \to A^\times$ and $S \in \text{Fin}(N) \to \text{Matrix}(m, m, \mathbb{R})$, where $A$ is a commutative ring. Suppose $S$ realizes $s$ and the scale is fully degenerate (isotropic): $s_a = s_b$ for all $a, b \in \text{Fin}(N)$. Then for any two directions $a, b \in \text{Fin}(N)$, the bracket of their scaling operators vanishes:
\[
\text{br}(S_a, S_b) = 0.
\]
\begin{proof}
Since the scale is fully degenerate, $s_a = s_b$ for all $a, b$. Because $S$ realizes $s$, we have $S_a = S_b$. Thus $\text{br}(S_a, S_b) = [S_a, S_b] = [S_a, S_a] = 0$ by the self-bracket property.
\end{proof}
\end{theorem}

\begin{theorem}[\lean{Dynamics.multiplets\_fixed\_by\_pattern}]
Let $s, t \in \text{Fin}(N) \to A^\times$ be two scale patterns. If they have the same coincidence structure,
\[
s_a = s_b \iff t_a = t_b \quad \text{for all } a, b \in \text{Fin}(N),
\]
then the multiplet structures determined by these patterns are identical: for all $a, b \in \text{Fin}(N)$,
\[
\text{SameOrbit}_s(a, b) \iff \text{SameOrbit}_t(a, b).
\]
\begin{proof}
The multiplets are defined by the equivalence relation on directions given by their scale values. Two patterns with identical coincidence structure define the same equivalence relation, hence the same multiplet structure.
\end{proof}
\end{theorem}

\begin{theorem}[\lean{Dynamics.multiplet\_change\_needs\_pattern\_change}]
Let $s, t \in \text{Fin}(N) \to A^\times$ be two scale patterns. If the multiplet structure changes between them,
\[
\exists a, b \in \text{Fin}(N) : \neg(\text{SameOrbit}_s(a, b) \iff \text{SameOrbit}_t(a, b)),
\]
then their coincidence structures must differ:
\[
\exists a, b \in \text{Fin}(N) : \neg((s_a = s_b) \iff (t_a = t_b)).
\]
\begin{proof}
Suppose for contradiction that the coincidence structures are identical: $(s_a = s_b) \iff (t_a = t_b)$ for all $a, b$. Then by the multiplets fixed by pattern theorem, the multiplet structures are identical, contradicting the assumption that they differ. Therefore the coincidence structures must differ.
\end{proof}
\end{theorem}

